%% Bilan énergétique d'un convertisseur (diagramme sagittal) :
%% énergie absorbée -> convertisseur -> énergie utile, plus les pertes.
\documentclass[tikz,border=4pt]{standalone}
\usepackage{liaisons}
\usetikzlibrary{shapes.geometric, positioning}
\begin{document}
\begin{tikzpicture}[x=1cm,y=1cm,
  bloc/.style={draw, line width=1pt, rounded corners=3pt, align=center,
               text width=1.5cm, minimum height=1cm, inner sep=2pt,
               font=\footnotesize\bfseries},
  conv/.style={ellipse, draw=solideE, text=solideE, line width=1pt,
               minimum width=2.3cm, minimum height=1.1cm, inner sep=1pt,
               font=\footnotesize\bfseries},
  grosse/.style={-{Stealth[length=8pt,width=7pt]}, line width=2.2pt}]

%% ---- chaîne principale ---------------------------------------------------
\node[bloc, draw=solideD, text=solideD] (abs) at (0,0)   {Énergie absorbée};
\node[conv]                             (cv)  at (3.1,0) {Convertisseur};
\node[bloc, draw=solideB, text=solideB] (ut)  at (6.1,0)  {Énergie utile};

\draw[grosse, solideD] (abs) -- (cv);
\draw[grosse, solideE] (cv)  -- (ut);

%% ---- pertes --------------------------------------------------------------
\node[bloc, draw=solideB!70, text=solideB!70, text width=1.1cm,
      minimum height=0.7cm] (pertes) at (3.1,-1.9) {Pertes};
\draw[grosse, solideE] (cv) -- (pertes);

%% ---- rappel du rendement -------------------------------------------------
\node[font=\small, align=center, anchor=north] at (3.1,-2.65)
  {$\eta = \dfrac{E_{\text{utile}}}{E_{\text{reçue}}}$\quad
   \footnotesize toujours entre 0 et 100~\%};
\end{tikzpicture}
\end{document}
